\chapter{Combat, Tactics, and Party Management}
\label{written-combat}

Three active companions give every action a noticeable place in a battle.
A timely heal can keep an attacker on its feet, a group spell can clear a
dangerous cluster, and a defensive action can buy the party the round it
needs to recover. The available companions and their equipment give the
player several ways to handle the same encounter.

Watch the enemy groups, the party's condition, and the resources needed
for the rest of the journey. A comfortable victory leaves enough HP, MP,
and useful items to continue exploring. Against a stronger opponent, the
same habits help the party settle into a sustainable sequence of attacks,
support, and recovery.

\section{Battle commands, target groups, turn order, and escape}

\subsection*{The opening battle menu}

The first battle menu offers four main actions:
\begin{description}
  \item[Fight:] Begin choosing the party's actions for the round. The
    current tactics determine which companions receive direct orders.
  \item[Change:] Adjust the active members when reserve companions and
    the encounter's conditions allow it.
  \item[Tactics:] Choose the battle behavior of computer-controlled
    companions or use the direct-command setting.
  \item[Run:] Attempt to leave an encounter that permits escape.
\end{description}

For a companion receiving direct orders, the available commands include
Attack, Spell, Parry, and Item. Attack uses the equipped weapon. Spell
opens the companion's battle abilities. Parry is the defensive action,
reducing incoming damage during that turn. Item opens the objects carried
by that companion.

Select the action and its target, then continue through the other directly
controlled members. Check the party's HP while choosing the round. A
companion who is already badly injured may need healing or a defensive
action before taking another heavy hit.

\subsection*{Reading the enemy groups}

The enemy display identifies the species present and the number remaining
in each group. The target menu uses those entries when selecting an
offensive action. A single-target attack affects one creature in the
chosen group, a group attack affects that group, and an all-enemy action
reaches across the opposing formation.

Choose targets according to their immediate threat. An enemy that heals,
calls reinforcements, drains MP, or inflicts a disabling condition can be
worth removing early. Concentrating several attacks on its group helps
reduce the number of dangerous actions the party faces in later rounds.

Group attacks are useful when several creatures of the same kind are
present. An all-enemy attack becomes useful when the opposing formation
contains several groups. Check which enemies remain as each new round
begins and choose the next targets accordingly.

\subsection*{Action order and recovery timing}

Agility influences when companions and enemies act. The order can vary
between rounds, so allow for an enemy reaching an injured ally before the
planned heal. A fast companion carrying a Medical Herb can provide another
recovery option when the main healer is slower or needs a different action.

At the start of a difficult round, identify the most urgent job. This may
be restoring a vulnerable ally, curing a disabling condition, reducing
incoming damage, or finishing a dangerous enemy. Assign the remaining
members around that action.

Parry gives a threatened companion a defensive turn while another member
handles recovery. After the party is healthy enough to resume its attack,
return that companion to its usual role. Keep enough recovery available
to handle the next unfavorable sequence of enemy actions.

\subsection*{Choosing to escape}

Run is useful when an encounter threatens the party's remaining supplies
or the journey has already become too expensive in HP and MP. A successful
escape lets the party continue from the field.

A failed attempt gives the enemies an opportunity to act. Consider the
party's condition before attempting another escape, especially when one
more enemy attack could bring down a companion. Healing or protecting the
party can make the next attempt safer.

Some fixed battles and boss encounters restrict escape. Prepare for those
encounters with a recovered party, accessible items, and a clear plan for
the first few rounds.

\section{Physical attacks, critical hits, multi-target weapons, and damage}

\subsection*{Attack power and the enemy's defense}

An ordinary physical attack uses the companion's strength and equipped
weapon against the enemy's physical defense. The resulting damage varies
between strikes. A stronger weapon improves the attacker's opportunities,
and lowering the enemy's defense can make the whole party's physical
attacks more effective.

Watch the damage produced by the current equipment. A sturdy enemy may
be easier to handle after Sap or Decrease has reduced its defense. Bikill
strengthens an ally's physical offense and is useful when that ally will
spend several rounds attacking.

For a short encounter, attacking immediately may finish the group quickly.
In a longer battle, a successful support action has time to improve several
later attacks. Choose the setup according to the opponent's durability
and the party's condition.

\subsection*{Choosing an attack pattern}

Weapon properties determine how far an attack reaches:
\begin{description}
  \item[Single-target weapons:] Ordinary swords, spears, claws, and similar
    weapons concentrate their attack on one creature in the selected group.
    They suit focused attacks against a dangerous enemy or a boss.
  \item[Group weapons:] Whips and weapons such as the Morning Star attack
    grouped enemies. They are useful when several creatures share a target
    entry. The Gringham Whip is a powerful example of this pattern.
  \item[All-enemy weapons:] Boomerangs, including the Bladed Boomerang and
    Flame Boomerang, reach across the enemy formation. The Iron Ball of
    Decimation also attacks all enemies.
  \item[Repeated attacks:] The Falcon Sword and Killer Earring produce
    two attacks from an attack action. Their usefulness depends on the
    wielder's strength, the enemy's defense, and the damage of each strike.
\end{description}

A Boomerang's 15 attack power and a Bladed Boomerang's 25 attack power can
be useful early because their attacks reach the whole opposing formation.
Give a group-clearing weapon to a companion who frequently handles ordinary
encounters. Keep a suitable focused weapon available for a battle with
one durable opponent.

Equipment compatibility remains important. Match the weapon's attack
pattern to a companion who can equip it and use it effectively. After a
change, check the resulting attack value and observe the next few battles.

\subsection*{Critical hits and special weapon effects}

Critical hits produce heavy physical damage and are especially valuable
against enemies with very high defense. They can also make an enemy's
ordinary attack suddenly dangerous. Keep enough HP available to survive
an unfavorable hit when exploring an unfamiliar area.

Some weapons have additional effects. The Poison Needle can produce a
one-hit kill, the Demon Spear also has a possible one-hit-kill effect,
and the Demon Claw can poison an enemy. Their special properties deserve
consideration alongside their attack power and eligible users.

The Dragon Killer is useful against dragons. Other equipment may add an
elemental property or an item-use action. Choose those properties for
the opponents the party is actually fighting, and review the result when
moving to another area.

\section{Spells, special abilities, MP, and reusable battle items}

\subsection*{Matching a spell to the formation}

The Spell command includes the companion's available battle spells and
abilities. Check the target pattern before choosing the action. Blaze
costs 2 MP and targets one enemy, Firebal costs 4 MP and targets a group,
and Bang costs 5 MP and attacks all enemies.

Use a single-target spell to concentrate damage on one important opponent.
A group spell is useful against a cluster of the same species. An all-enemy
spell helps when several groups are present. Enemy resistances affect the
result, so watch how well the selected spell works in the current encounter.

A well-chosen offensive spell can shorten a battle and reduce the healing
needed afterward. Keep the journey's remaining demands in mind when
deciding how much MP to spend on an ordinary group.

\subsection*{Monster abilities}

Monster companions bring abilities such as breath attacks, dances,
physical techniques, and disruption effects. Breath attacks and many
monster-specific actions provide useful options with no MP cost.
Their target patterns and effects still matter when choosing an action.

A breath user can contribute consistent damage through a long dungeon.
A companion with sleep, paralysis, or disruptive abilities can help control
a dangerous group. Inspect the available abilities after recruitment and
after level gains so that the new options become part of the party's plan.

An ability's usefulness depends on the opponent. Switch to another useful
action when a creature resists the chosen effect. Physical damage, another
spell family, support, or an item can give that companion a productive turn.

\subsection*{Healing spells and their costs}

Heal costs 2 MP, Healmore costs 5 MP, and Healall costs 7 MP. Healall restores
one ally's HP completely. Choose a recovery spell that suits the injury
and leaves enough resources for the next rounds.

Healus costs 18 MP and restores HP across the active party. Healusall costs
36 MP and provides full-party HP recovery. These spells become valuable
when an enemy damages several companions at once.

Revival and status treatment add to the healer's workload. Keep another
recovery option available so that the party can handle an injury while
the main healer deals with a fallen or disabled companion.

\subsection*{Reusable item actions}

Certain possessions provide battle effects through the Item command:
\begin{description}
  \item[Staff of Benediction:] Provides a Healmore effect when used in
    battle. It gives its eligible user an additional recovery action.
  \item[Staff of Antimagic:] Provides a Stopspell effect, useful against
    an enemy that depends on spells and is vulnerable to the effect.
  \item[Sage's Stone:] Restores HP to the active party and can be used
    repeatedly during battle.
  \item[Drum of War:] Strengthens the active party's physical offense,
    making it useful when several members will continue attacking.
\end{description}

Give a useful reusable object to an available companion and select its
item action during that companion's turn. Its use occupies an action,
so include that action in the round's plan. A fast user can be useful
when the effect needs to arrive early.

Reusable recovery can preserve a caster's MP for revival, strong healing,
or support. Keep consumables available for urgent situations where another
companion needs to handle the same job.

\section{Status ailments, buffs, debuffs, resistance, and spell reflection}

\subsection*{Responding to disabling effects}

Poison, paralysis, sleep, confusion, and spell sealing can change the
party's priorities in a single round. Check which companions can still
act and which recovery tools are available to them.

Antidote and Antidote items treat poison. Numboff and Full Moon Herbs
treat paralysis. An Angel's Bell cures confusion, has unlimited uses,
and works during battle.
Carry the relevant items with active members who can use them when an
important caster is disabled.

Sleep and confusion can make a planned action unreliable. Protect the
remaining available companions, keep recovery ready, and remove the
enemy causing the condition when practical. If spell sealing affects
the healer, use accessible item recovery while handling the encounter.

Some enemy effects drain MP. Removing that enemy early can preserve the
resources needed for the rest of the dungeon. Check the healer's remaining
MP after the battle before continuing.

\subsection*{Strengthening the party and weakening an enemy}

Upper raises one ally's physical defense, and Increase strengthens the
party's physical defense. They are useful against an opponent that relies
on ordinary physical attacks. Bikill improves an ally's physical offense.
Choose a companion likely to keep attacking for several rounds when using it.

Sap lowers a single enemy's defense, and Decrease works against a group.
A successful defense reduction can improve the results of several party
members' attacks. A resistant opponent may require a different plan.

Some enemies remove beneficial effects during battle. Watch for that
action and reassess which support effects are worth restoring. Keep
healing and immediate survival ahead of rebuilding a full set of boosts
when the party has just taken heavy damage.

\subsection*{Breath protection and equipment resistance}

Barrier helps protect the active party from breath attacks. It is useful
when a strong fire or cold breath is a regular part of the opponent's
attack pattern. Combine that protection with equipment suited to the
same threat and enough party healing to cover the remaining damage.

Armor and accessories can provide protection from particular attacks or
status effects. Inspect these properties when preparing for an area.
A companion who repeatedly loses its turn to a condition may benefit
from a protective item as well as a ready cure.

Enemy resistance also affects offensive spells and disabling actions.
Watch the battle messages and damage results. A spell that works well
against one species can have a different result against another.

\subsection*{Using Bounce carefully}

Bounce reflects spells and can help protect a companion from a spell-heavy
opponent. Plan recovery for the protected member before relying on it.
Item-based healing provides a useful option while reflection is active.

An enemy using reflection also changes offensive spell selection.
Physical attacks and suitable abilities or item actions give the party
other ways to continue. Check the target's current condition before
committing a powerful spell.

\section{AI tactic modes and when to use direct commands}

\subsection*{Choosing a policy}

Tactics set the behavior of computer-controlled companions. The Hero's
own actions are chosen through the ordinary command menu. Choose the
policy that suits the current opponents, remaining resources, and length
of the journey.

The six policies cover the following approaches:
\begin{description}
  \item[All-out attack:] Prioritizes strong offensive actions and powerful
    available spells or abilities. It suits a group that the party wants
    to remove quickly. Watch MP and the need for a separate healing action.
  \item[Balanced fighting:] Mixes offense with recovery and useful support
    according to the situation. It is convenient for familiar ordinary
    encounters with a healthy, well-equipped party.
  \item[Support the Hero:] Emphasizes assistance for the Hero through
    recovery, strengthening, protection, and hindering the opposition.
    Use it with the Hero active, able to attack, and suitably equipped.
  \item[Conserve MP:] Uses physical attacks and available MP-free actions.
    It is useful during manageable encounters on the way to a demanding
    area. Keep an emergency recovery action available.
  \item[Protect life:] Places greater priority on healing and defensive
    support. It can help a party facing dangerous ordinary groups, and
    its frequent recovery can consume a substantial amount of MP.
  \item[Follow orders:] Gives direct control over eligible companions'
    actions. It is useful for precise targets, planned item use, and
    coordinated boss strategies.
\end{description}

Under an MP-conserving policy, the Hero's manually selected spells still
use their normal MP costs. Check the party's resources as the route
continues and change the policy when its priorities need to change.

\subsection*{When direct orders are useful}

Use Follow orders when a specific companion needs to heal, revive, use
a particular item, or apply a support effect at a planned moment. Direct
orders also help coordinate several attacks against a selected enemy group.

A reusable healing item has a particular owner and uses that owner's turn.
Selecting the action directly makes it easy to account for both. A second
member can handle single-target recovery while the item user restores the
active party.

For a boss encounter, decide the first few rounds in advance. Assign the
initial defense or offense boosts, identify the main recovery member,
and leave room for an urgent cure or revival. Adjust the plan as the
opponent reveals its pattern.

\section{Intelligence, obedience, and monster command restrictions}

\subsection*{Developing reliable companions}

A monster's intelligence affects its ability to follow commands
consistently. A low-intelligence recruit can act on its own or fail to
carry out the intended action. Keep a reliable recovery member available
when training such a companion.

Raising intelligence above 20 provides a practical target for reliable
obedience. Check the Status screen as the monster gains levels and review
its available spells and equipment at the same time.

Intelligence Seeds provide permanent increases. Use them on the intended
companion and check the resulting value after each use. This is useful
for a promising species whose natural intelligence growth or level cap
leaves its command reliability short of what the party needs.

Bianca's Ribbon also raises intelligence when equipped by an eligible
monster. Check the resulting Status value when assigning it and keep
the Ribbon available for the later journeys.

\subsection*{Level caps and a companion's role}

Maximum level varies by species. The Dancing Jewel reaches its cap at
level 7, the Orc King at 30, and the Curer at 50. A Slime Knight can continue
to level 99. Those limits give the companions different growth schedules.

Consider the complete set of abilities available at the current level.
A companion with a low cap can still have useful support or protection,
and a longer-growing species may need several levels before an important
healing spell becomes available.

Keep an unreliable new recruit in a suitable training role while another
member handles essential recovery. Once its intelligence and abilities
support the intended job, move it into a position with more responsibility.

\subsection*{Temporary companions}

Temporary companions have different battle behavior. Papas fights
autonomously. Bella provides healing and support during the faerie
journey. Check the selected tactic and each companion's available
commands when planning the Hero's actions around their assistance.

Review the lineup after a temporary companion joins or leaves. Check
the available healing, item owners, and equipment before continuing into
the next difficult area.

\section{The three-member active party and wagon reserves}

\subsection*{Filling the active positions}

The normal active party has three members. With the wagon available,
the travelling roster can contain eight members, leaving five reserve
positions behind the active three.

Choose a lineup whose abilities work together. Useful jobs include
physical damage, recovery, support, and attacks that reach several enemies.
One companion can cover more than one job, depending on its spells,
equipment, and current level.

For ordinary exploration, a group-clearing attacker, a reliable recovery
member, and a durable companion make a comfortable combination. Before
a stronger opponent, review its likely threats and adjust the roles.
Several physical attackers can benefit from offensive support, while
frequent party-wide damage creates more work for recovery.

\subsection*{Companions with useful combinations}

A Slime Knight combines physical equipment options with healing. It
starts with Heal, learns Healmore at level 5, and gains Healall at level
20. Robmagic can draw MP from susceptible enemies. Antidote at level 9
and Numboff at level 18 add poison and paralysis treatment. These tools
give it useful recovery work alongside its attacks.

An Orc King starts with Vivify, learns Healus at level 8, and gains Revive
at level 11. Its developing spell list can cover both revival and
party-wide healing. Keep enough MP available for those jobs when relying
on it during a longer encounter.

A Curer begins with Healall and later learns additional recovery and
support. Healmore at level 3 provides a cheaper single-target recovery
choice, and Upper at level 8 protects one ally. It gains Increase at
level 10, Vivify at 12, Healus at 18, Revive at 23, and Healusall at 30.
Its role can expand substantially as it grows.

Use these combinations when distributing the party's responsibilities.
A companion capable of fighting and healing can cover an urgent injury
while the usual healer revives another member or restores the group.

\subsection*{Making the reserves useful}

Reserve positions can hold another healer, an alternative attacker,
a companion with a particular support ability, or a recruit being trained.
Keep their equipment and condition suitable for the job they may need
to perform.

Where the wagon is available, changing members can bring fresh MP or
a more useful ability into the active lineup. A reserve holding an
important battle item must be brought into an available position to use
it during its turn.

Review both the active party and reserves before leaving town. Healing
and supplies distributed across the travelling roster give more options
during a long sequence of battles.

\section{Changing members, wagon access, experience, and dungeon restrictions}

\subsection*{Changing the lineup}

Use the battle menu's Change command when the wagon and encounter permit
it. Select the available companions needed for the next part of the fight,
then review the new lineup before choosing its actions.

A fresh healer can take over when the original member has spent most of
its MP. A durable attacker can help after a vulnerable companion has
taken heavy damage. A specialist can enter to supply a protective spell
or a particular item action.

The incoming companion needs its own equipment, available MP, and useful
possessions. Prepare these before the fight so that the change provides
an immediate benefit.

\subsection*{Areas with restricted wagon access}

Some areas admit only the active party. Choose the three companions needed
for the dungeon before entering, and give them the supplies they will
need inside. Include reliable healing and a way to handle the area's
important status threats.

Check access when the route moves through an entrance, floor connection,
or fixed encounter. The companions available in an open-field battle can
differ from the members available deeper in a dungeon.

Keep a useful battle item with a companion who will actually enter the
area. The same preparation applies to required keys, recovery supplies,
and any field spell needed for the way out.

\subsection*{Experience and training}

Victories award experience to eligible living companions. In battles
where the wagon is available, reserves can gain experience alongside the
active members. This gives newly recruited companions a way to develop
while a dependable front line handles the encounters.

Choose manageable groups when training a fragile recruit. Keep it alive
and recovered, and inspect its new abilities as levels increase. A useful
healing or support spell may arrive well before its final level cap.

For a dungeon with restricted access, arrange the party around the members
participating in that area. Use accessible wagon battles when the goal
is to develop companions waiting in reserve.

\section{Healing economy, resource conservation, and boss preparation}

\subsection*{Keeping a dungeon run sustainable}

The journey uses HP, MP, and consumables across several encounters. Keep
the main healer's remaining MP visible when deciding how far to explore.
Medical Herbs can handle smaller injuries and preserve spells for the
next difficult stretch.

MP-free abilities and reusable items can provide steady work through a
long route. Assign those actions to companions who can use them reliably,
and keep a stronger recovery option available for a dangerous round.

When escape spells are part of the return plan, reserve their cost.
Outside and Return each cost 8 MP. A companion casting both needs 16 MP
available for that journey home, plus any healing still needed beforehand.

\subsection*{Preparing for a boss}

Rest the party, treat ailments, and check equipment before the approach.
Fill the available inventories with useful recovery and status-treatment
items. Arrange item ownership so that several active members can help
when the regular healer is occupied.

Use the opponent's attack pattern to choose support. Increase helps with
physical pressure, Barrier helps with breath attacks, and suitable armor
can provide further protection. Offensive strengthening becomes useful
when the party can keep an attacker working for several rounds.

Some enemies perform several actions in a round. Allow enough HP for
the possible sequence of attacks before relying on the next heal. Keep
a second recovery action available for a round in which one companion
must be revived or cured.

\subsection*{A workable rhythm for a long fight}

Begin by securing the party's immediate safety and applying useful
support. Continue with focused attacks while watching HP and MP. A
single-target healer can handle an isolated heavy hit, and party-wide
recovery can cover a breath or spell that injures all three members.

The Sage's Stone is useful for repeated party recovery when it is
available. Its user spends an action restoring the group while another
companion handles a more urgent single-target heal, revival, or support
spell.

The Drum of War supports a physically focused active party when available.
After an enemy removes strengthening effects, restore the most useful
ones as the party's condition allows. Keep the recovery plan working
through those changes.

If a healer is running low on MP and the wagon is accessible, bring in
a prepared reserve. Check the incoming member's item ownership and
abilities before relying on the next round.

\section{Enemy summoning, metal enemies, defeat, and recovery}

\subsection*{Handling reinforcements}

Some enemies call additional creatures into the battle. Reinforcements
increase the number of actions the party must handle and can prolong
an encounter that initially looked manageable.

Remove the summoner early when the party can do so safely. Keep track
of the remaining groups and the abilities of newly arrived creatures.
A called healer can make an enemy group harder to finish, while another
disabling enemy can increase the need for recovery items.

Called creatures can also matter during recruitment. Healing monsters
can appear through summoning, so watch for the species being pursued
and manage the battle with that goal in mind. Keep the party healthy
enough to finish the encounter.

\subsection*{Fighting metal enemies}

Metal enemies offer valuable experience and can flee quickly. Their high
defense and strong resistance to offensive magic make physical actions
and critical hits important parts of the plan.

Use direct commands to concentrate suitable attacks on the metal group.
A fast companion has a useful chance to act before the enemy escapes.
Repeated-attack weapons such as the Falcon Sword or Killer Earring can
provide additional physical strikes when an eligible user has them.

Keep ordinary enemies in the encounter under control as well. A harmful
status effect or a badly injured companion can interrupt the attempt.
After a successful battle, check the levels and abilities of companions
who gained experience, including eligible wagon reserves.

Choose a hunting area the party can handle comfortably. Nearby recovery
and dependable equipment help make repeated attempts worthwhile, and
ordinary victories along the route still contribute to development.

\subsection*{Recovering after a difficult encounter}

Treat poison and other persistent conditions, restore enough HP for the
next stretch, and check the remaining MP and items. Reorganize the active
party when another companion is better prepared to continue.

If the party is defeated, half of the carried gold is lost. Banked gold
remains protected. Arrange revival, treatment, and rest after returning
to the recovery point, then reconsider the equipment and supplies for
the next attempt.

Before returning to the difficult area, identify what caused the collapse.
Repeated breath damage may call for protection and group recovery.
Spell sealing makes accessible healing items valuable. A sequence of
heavy physical hits may require better defense, a sturdier companion,
or earlier healing.

Record the recovered adventure at a church once the party is ready.
The next attempt can then begin with its equipment, recovery roles,
and travel plan already in place.

